\chapter{Nuclear Medicine Cardiac Imaging}

\section*{What Is This Test?}
Nuclear cardiac stress tests use radioactive tracers (thallium, sestamibi) to show blood flow to the heart muscle at rest and stress. Areas with reduced flow during stress but normal flow at rest indicate coronary artery disease.

\section*{Alternative Names}
Nuclear stress test, myocardial perfusion imaging (MPI), cardiac PET, SPECT cardiac imaging, nuclear cardiology, sestamibi scan, thallium stress test


\section*{Principle}
Myocardial perfusion imaging uses Tc-99m sestamibi or tetrofosmin injected at stress and rest. The tracer distributes proportionally to myocardial blood flow. SPECT imaging detects perfusion defects: reversible defects indicate ischemia; fixed defects indicate infarct. ECG-gating evaluates wall motion and ejection fraction.

\section*{History}
Nuclear cardiology began with the use of thallium-201 for myocardial perfusion imaging in the 1970s. The development of technetium-99m labeled agents in the 1980s improved image quality and reduced radiation exposure.

\section*{Why This Test Is Ordered}
Your doctor orders a nuclear stress test (also called myocardial perfusion imaging) when you have chest pain, shortness of breath on exertion, or other symptoms suggesting your heart is not getting enough blood. It answers the question: is your coronary artery blood flow adequate at rest and during stress, and if not, which areas of your heart are affected?

\section*{Contraindications and Precautions}
Nuclear cardiac imaging uses a small amount of radioactive tracer and is generally not performed during pregnancy. Inform the team if you are pregnant or breastfeeding. The tracer passes into breast milk; discuss timing with your doctor.

\section*{Risks}
The radioactive tracer exposes you to a small amount of radiation, comparable to a CT scan. Rare allergic reactions to the tracer may occur. The exercise or pharmacological stress portion carries the same risks as any stress test (chest pain, shortness of breath, arrhythmia). Inform the team if you are pregnant or breastfeeding.

\section*{What Happens After the Test}
You can resume normal activities. Results typically take one to three hours for immediate interpretation and one to three days for the formal report.

\section*{Understanding Your Results}

\section*{Normal Results}
A normal result shows uniform blood flow to the heart muscle at rest and during stress, with no reversible or fixed perfusion defects.

\section*{Follow-up Tests}
If the nuclear stress test is abnormal, your doctor may order coronary angiography to visualize blockages directly, a cardiac MRI for tissue characterization, an echocardiogram to assess heart function, or a CT coronary angiogram for non-invasive evaluation of coronary arteries.
\clearpage
